\documentclass[11pt]{article}
\usepackage{fontspec}
\setmainfont{Times New Roman}
\setsansfont{Arial}
\setmonofont{Consolas}
\usepackage{amsmath,amssymb}
\usepackage{geometry}
\usepackage{microtype}
\geometry{a4paper,margin=3cm}
\setlength{\parindent}{1.25em}
\setlength{\parskip}{0pt}
\emergencystretch=6em
\allowdisplaybreaks

\begin{document}

\begin{center}
{\Large\bfseries 4. До інваріантної теорії форм від $n$ змінних}\par
\vspace{1.2em}
\emph{Journal f. d. reine u. angew. Math.} 139 (1911), с. 118--154
\end{center}

\section*{Вступ.}
У проективній інваріантній теорії форм від $n$ змінних питання, що прилягають до головної проблеми -- доведення скінченності системи форм, -- майже не розглядалися, якщо вийти за межі бінарної та тернарної областей. Це питання про взаємну залежність форм і про методи опанування цього зв'язку між формами. Ці методи істотно спираються -- як це повністю доведено в бінарній і тернарній областях -- на дві теореми, які Study назвав ``фундаментальними теоремами символічного методу'', а саме на теорему про символічну зображуваність форм і теорему про символічні тотожності.\footnote{E. Study, \emph{Methoden zur Theorie der ternären Formen} (Leipzig, Teubner, 1889). II. \S\ 6. -- Про появу цих фундаментальних теорем також в інваріантній теорії спеціальних груп перетворень пор. Study, \emph{Das Apollonische Problem} (\emph{Math. Ann.} Bd. 49 [1897]) і \emph{Zur Differentialgeometrie der analytischen Kurven} (\emph{Transactions of the American Math. Society}, Bd. 10 [1909]; в оригінальному виданні помилково надруковано Bd. 1).} Остання теорема стверджує, що всі відношення між цілими інваріантними утвореннями одержуються послідовним застосуванням невеликої кількості символічних тотожностей; отже, символічний метод, і тільки він, дає знання зв'язку між формами, не виходячи з області інваріантів.

На причину труднощів під час переходу до $n$ змінних уже вказав Gordan у своїй Ерлангенській програмі;\footnote{P. Gordan, \emph{Über das Formensystem binärer Formen}. S. 46. (Leipzig, Teubner, 1875).} вона полягає в тому, що перша теорема символічного методу в її загальному формулюванні вже не справджується. Як відомо, в $n$-арній області з'являються рядки змінних вищого рівня $p_\rho$\footnote{Щодо позначення пор. \S\ 1.} і, аналогічно, символьні рядки вищого рівня $S^\rho$; до цих рядків приходять як тоді, коли, за первісним підходом Clebsch,\footnote{A. Clebsch, \emph{über eine Fundamentalaufgabe der Invariantentheorie}. (\emph{Abh. der Ges. d. Wiss. zu Göttingen}, Bd. XVII, 1872).} виходять із форм із довільно багатьма рядками $p_\rho$, так і тоді, коли виходять із форм, що містять лише довільно багато когредієнтних рядків $x$, за підходом F. Deruyts і A. Capelli.\footnote{Пор. A. Capelli, \emph{Lezioni sulla teoria delle forme algebriche} (Napoli, 1902), \S\S\ 22--24, і наведену там літературу.} Символічного подання, яке явно\footnote{Точніше означення ``явного подання'' у \S\ 2.} містило б рядки $p_\rho$, $S^\rho$, тепер не існує; потрібно розкласти $p_\rho$, $S^\rho$ на окремі рядки $x$ і контрагредієнтні до них $s$ та розподілити їх між різними детермінантами. Щоправда, за допомогою диференціальних умов (пор. формулу (19.)) можна перевірити, що даний агрегат детермінантів має властивість залежати тільки від рядків $p_\rho$, $S^\rho$; але таким способом не можна здобути огляду всієї сукупності інваріантних утворень і процесів їх породження.\footnote{Також обчислювально дуже цінний дальший розвиток символіки, що походить від Study (повідомлений F. Engel, \emph{Ber. d. K. Sächs. Ges. d. Wiss., math.-phys. Kl.}, 1900, S. 126, а для кватернарної області -- E. Study, \emph{Geometrie der Dynamen}, S. 135), не є явним поданням у нашому розумінні. Наприклад, він навряд чи привів би до несимволічних процесів диференціювання для породження форм.}

Тепер ми покажемо в \S\S\ 2 і 6, як, застосовуючи ``теорему про відповідні матриці'',\footnote{Якщо не рахувати Graßmann-ової \emph{Ausdehnungslehre} 1862 року, Nr. 112, це вперше з'являється в A. Brill, \emph{über zwei Berührungsprobleme}, \S\ 2 (\emph{Math. Ann.} Bd. 4, 1871).} можна знову одержати першу фундаментальну теорему в її загальності, замінивши подання через детермінанти поданням через ``матричні добутки''. У \S\ 2 показано, що кожен матричний добуток, який явно містить рядки $S^\rho,p_\rho$, є інваріантним утворенням основної форми; обернення, наведене в \S\ 6, а саме що кожне інваріантне утворення можна явно подати матричними добутками, вдається тільки через перенесення другої фундаментальної теореми (\S\S\ 3--5).

Ця теорема відома для форм, що містять лише рядки змінних $x$.\footnote{E. Pascal, \emph{Giorn. di mat.} 26 (1888), S. 33, 102; \emph{Rendic. d. Accad. d. Linc.} (4) 4 (1888), S. 119; ib. \emph{Mem.} (4) 5 (1888), S. 375.} Загальну проблему -- встановити всю сукупність нерозкладних тотожностей, у яких усі рядки $S^\rho,p_\rho$ розглядаються як нерозкладні, -- сформулював Study;\footnote{\emph{Methoden \dots}, S. 204, примітка 17).} водночас він бачить у кількості тотожностей, що виникають, міру трудності методів у $n$-арній області. Оскільки тепер вдається об'єднати всю сукупність тотожностей в одну формулу (36.), цю трудність принаймні зведено до мінімуму. У застосуваннях, однак, часто з'являтимуться певні вирізнені спеціальні випадки (36.), як-от (20.), (21.), що характеризуються тим, що допускають явне подання без підстановки нових рядків; або формула (31.), яку названо ``тотожністю розкладання'', бо рядок $p_\rho$ у ній ніби розкладено на рядки $y$.

На двох фундаментальних теоремах тепер легко побудувати дальшу теорію. Перша теорема безпосередньо дає процеси породження форм, символічний процес згортання і несимволічні процеси диференціювання,\footnote{У праці \emph{Invariante Gebilde von Nullsystemen} (\emph{Sitzungsber. d. Ak. d. Wiss. Wien}, Bd. 97, Abt. IIa, 1888) Mertens іншим, не безпосередньо узагальнюваним шляхом прийшов до кватернарного процесу згортання. Альтернувальний інваріант від 6 рядків $S^2$, що там виникає, відрізняється від нашого, який виникає при одноразовому застосуванні підстановки (11.), агрегатом парних інваріантів. Для одного спеціального випадку кватернарний процес згортання є також у P. Gordan, \emph{Das simultane System von zwei quadratischen quaternären Formen} (\emph{Math. Ann.} Bd. 56, 1903).} тоді як тотожність розкладання усуває серед цих процесів усі еквівалентні (\S\ 6). Знання процесів диференціювання далі дає перенесення поняття ``нормальної форми'' на $n$-арну область і, за допомогою способу доведення, наведеного Mertens\footnote{F. Mertens, \emph{Über invariante Gebilde quaternärer Formen}. (\emph{Sitzber. d. Ak. d. Wiss. Wien}. Bd. 98, Abt. IIa, 1889.)} у кватернарній області, теорему про те, що систему інваріантних утворень можна замінити еквівалентною системою нормальних форм (\S\ 7). За цієї останньої передумови я показала у своїй дисертації,\footnote{\emph{Über die Bildung des Formensystems der ternären biquadratischen Form}. (Цей журнал, Bd. 134, 1908.)} що тернарний процес згортання можна породити послідовним застосуванням двох основних згортань. На підставі цієї теореми і теорії розкладів у ряди я потім, запровадивши поняття ``ряду форм'', розвинула головні теореми редукції для систем форм. Цілком аналогічно в $n$-арній області процес згортання можна породити з $n-1$ основних згортань (\S\ 8), і можна встановити теореми редукції (\S\ 9).

\paragraph{Післямова.} З приводу Зальцбурзького повідомлення проф. Е. Мюллер із Відня звернув мою увагу на те, що обчислення з матричними добутками, покладене в основу цієї праці, за принципом тотожне з численням Graßmann-ової теорії протяжності.\footnote{\emph{Hermann Graßmanns gesammelte mathematische und physikalische Werke. Ersten Bandes zweiter Teil: Die Ausdehnungslehre von 1862. In Gemeinschaft mit H. Graßmann d. Jüngeren herausgegeben v. Fr. Engel} (Leipzig, Teubner, 1896).} Справді, Graßmann-ів ``комбінаторний добуток'' $[S^{\rho_1}S^{\rho_2}\cdots S^{\rho_i}]$ можна розглядати як матрицю, задану цими рядками (пор. \S\ 2), причому ``прогресивний добуток'' є матрицею самих рядків, ``регресивний'' -- матрицею зіставлених рядків $S_{n-\rho_1},S_{n-\rho_2},\ldots,S_{n-\rho_i}$, а матриця, що відповідає ``мішаному добутку'', виникає тоді, коли кілька рядків $S$ об'єднуються підстановкою (9.) в нові рядки $P,Q$. Наш ``зіставлений рядок'' відповідає при цьому Graßmann-овому ``доповненню'', а наш ``матричний добуток'' -- ``мішаному добутку ступеня 0''.

За такого зіставлення розклад, наведений в \emph{Ausdehnungslehre} 1862 року, Nr. 113, стає тотожним із формулою, дуалістичною до нашої формули (32.). У спеціальному випадку $\rho=1$ формула (32.) переходить у (17.), а дуалістична формула -- у (16.). Із цих спеціальних випадків Graßmann-ового розкладу Мюллер нещодавно вивів згадані вище тотожності (20.), (21.).\footnote{E. Müller, \emph{Beiträge zur Graßmannschen Ausdehnungslehre. 1. Mitteilung: Einige allgemeine Sätze}. (\emph{Sitzungsber. d. Ak. d. Wiss. Wien}; Bd. 118, Abt. IIa, Juli 1909), Satz 16 і 18.}

На відміну від Graßmann--Müller-ового виведення, наше виведення водночас дає доведення повноти тотожностей, що для питань, які тут розглядаються, видається найістотнішим; від (20.), (21.) воно далі переходить до найзагальніших нерозкладних тотожностей (36.), які, здається, досі ще не були помічені.

\end{document}
